\section{Conclusion}
\label{sec:conclusion}

This paper evaluated explanation-guided active learning and controlled TreeSHAP
diagnostics for static Windows PE malware classification. The acquisition family
combines predictive uncertainty with directional entropy and signed-attribution
variance. On an identity-audited 2,000-sample EMBER subset, evaluated with five
outer folds and five seeds per fold, $U+H$ achieved the highest ROC/AP learning-
curve area. Its improvements over Random were positive in every outer fold;
improvements over Uncertainty were smaller and occurred in four of five folds.
$U+V$ reached the predeclared validation target with the fewest mean training
labels, while $\xfull$ reduced the mean budget by 92 labels relative to Random.
No configuration was best across learning quality, target attainment, restricted-
FPR behavior, and computation.

The TreeSHAP study extended the evaluation to 25 fitted models. Top-ranked
coordinates produced substantially larger AUC decreases than equal-count random
coordinates under both training-mean and training-median replacement. Global and
top-k rankings showed moderate agreement and were less stable across data folds
than across seeds within one fold. These findings support relative reliance of
the fitted models on the ranked coordinates under the specified interventions.
They do not establish causal malware features or valid executable edits.

The results motivate a focused next experiment rather than a broader deployment
claim. Uncertainty plus diversity, calibrated committee acquisition, weight and
normalization sensitivity, stronger deletion controls, and external/temporal
evaluation are the primary evidence gates. Completing them would determine
whether the measured gains reflect uniquely useful explanation signals and
whether they transfer beyond the audited subset.

The study also illustrates a general reporting principle for security ML: label
efficiency, operating-point quality, attribution behavior, and reproducibility
answer different questions. Reporting their trade-offs and limitations provides
a stronger basis for model maintenance than selecting a single favorable AUC or
explanation image.

